\chapter{Implementation and Testing}
\label{ch:implementation}

This chapter details the implementation progress across all four iterations, the algorithm design, external APIs, testing results, and the full cross-framework benchmarking evaluation.

\section{Algorithm Design}

The core pipeline of QCanvas is its multi-stage quantum circuit compilation algorithm. The pseudocode below describes the main conversion flow.

\begin{algorithm}[H]
\caption{QCanvas Circuit Compilation Pipeline}
\begin{algorithmic}[1]
\Require Source code $S$, framework $F$, backend $B$, shots $n$
\Ensure Measurement counts $C$, execution metadata $M$

\State $F \gets \textsc{DetectFramework}(S)$  \Comment{Inspect imports and class patterns}
\State $\text{AST} \gets \textsc{ParseAST}(S, F)$  \Comment{AST traversal; runtime fallback if ambiguous}
\State $\text{ops} \gets \textsc{ExtractOps}(\text{AST})$  \Comment{Gate list, qubits, params, measurements}
\State $\text{QASM} \gets \textsc{QASM3Builder.build}(\text{ops})$  \Comment{Header + declarations + gates}
\State $\text{diag} \gets \textsc{Validate}(\text{QASM})$  \Comment{Subset check; emit line-precise errors}
\If{$\text{diag.has\_errors}$}
    \State \Return $\text{diag}$  \Comment{Return diagnostics to frontend; halt}
\EndIf
\State $\text{job} \gets \textsc{QSim.submit}(\text{QASM}, B, n)$  \Comment{Dispatch to selected backend}
\State $C, M \gets \textsc{QSim.collect}(\text{job})$  \Comment{Counts, exec time, state snapshot}
\State \Return $C, M$
\end{algorithmic}
\end{algorithm}

\section{External APIs and SDKs}

\begin{table}[!h]
\resizebox{\textwidth}{!}{%
\begin{tabular}{|p{3cm}|p{4cm}|p{4cm}|p{4cm}|}
\hline
\textbf{API / SDK} & \textbf{Description} & \textbf{Purpose} & \textbf{Key Interface} \\ \hline
Qiskit (v0.44+) & IBM quantum framework & Parse/convert Qiskit circuits to OpenQASM 3.0 & \texttt{QuantumCircuit} \cite{qiskit} \\ \hline
Cirq (v1.0+) & Google quantum framework & Parse/convert Cirq circuits & \texttt{cirq.Circuit} \cite{cirq_developers_2023} \\ \hline
PennyLane (v0.33+) & Xanadu variational framework & Parse/convert PennyLane QNodes & \texttt{qml.QNode} \cite{bergholm2018pennylane} \\ \hline
QSim SDK (PyPI) & Custom simulation kernel & Execute OpenQASM 3.0 circuits & \texttt{StatevectorBackend.run()} \cite{qsim_sdk_2026} \\ \hline
FastAPI (v0.110+) & Async Python web framework & REST and WebSocket API & \texttt{APIRouter}, \texttt{WebSocket} \\ \hline
Monaco Editor & Microsoft in-browser IDE & Syntax highlighting, hover providers & \texttt{editor.create()} \\ \hline
D3.js (v7) & SVG data-driven library & Real-time 2D circuit diagram & \texttt{d3.select}, \texttt{d3.join()} \\ \hline
PostgreSQL/SQLAlchemy & Relational DB + ORM & Persistent storage for all entities & \texttt{Session}, \texttt{Base.metadata} \\ \hline
\end{tabular}%
}
\caption{External APIs and SDKs Used}
\label{tab:external-apis}
\end{table}

\section{Iteration-by-Iteration Implementation Progress}

\subsection{Iteration 1: Core Infrastructure and Basic Features}
Iteration 1 established the foundational architecture: FastAPI backend, Next.js 14 frontend with TypeScript and Tailwind CSS, Docker Compose DevOps, and initial Cirq/Qiskit/PennyLane to OpenQASM 3.0 conversion with standard gates, if-else control flow, and the complete OpenQASM 3.0 type system \cite{cross2021openqasm}.

\subsection{Iteration 2: Advanced Features and Hybrid Execution}
Iteration 2 added advanced OpenQASM 3.0 constructs (\texttt{while}, \texttt{break}, \texttt{continue}, subroutines, complex types, bitwise operators), expanded PennyLane gate coverage (CY, CH, CRX, CRY, CRZ, CP, CSWAP, CCZ, GlobalPhase, \texttt{negctrl}, \texttt{ctrl(n)}, \texttt{pow(k)}), the full database layer (SQLAlchemy models, Alembic migrations, FK constraints), and the hybrid CPU--QPU execution model separating classical orchestration from QSim-powered quantum execution. Security was hardened with AES-256 key encryption, Bcrypt password hashing, and RBAC.

\subsection{Iteration 3: Educational Intelligence and Community Features}
Iteration 3 delivered the real-time D3.js circuit visualisation pipeline with 500ms debounce, the ``Quantum Explain It'' hover system, an event-driven gamification engine with 44 unique achievements and XP progression to Level~50+, and the community recipe-sharing platform. The plug-in language architecture was reinforced with Monarch tokenizer injection for \texttt{qasm} and \texttt{pythonq} language modes.

\subsection{Iteration 4: QSim SDK Integration, Benchmarking, and PyPI Deployment}
Iteration 4 marks the project's completion. Both the QCanvas converter SDK and the QSim simulation SDK were packaged and deployed to PyPI:
\begin{itemize}
    \item \texttt{pip install qsim-sdk} --- standalone quantum simulation engine.
    \item \texttt{pip install qcanvas-sdk} --- Qiskit/Cirq/PennyLane to OpenQASM 3.0 converter.
\end{itemize}
The example library was expanded to 40+ categorised algorithms. A comprehensive cross-framework benchmarking study was conducted and a Docker Compose production deployment with Nginx reverse proxy was finalised.

\section{User Interface Implementation}

The QCanvas Web IDE, delivered across all four iterations, provides:
\begin{itemize}
    \item Monaco editor with custom syntax highlighting and hover-based quantum explanations.
    \item Real-time D3.js 2D circuit diagram (debounced) and 3D visualisation via \texttt{@react-three/fiber}.
    \item Chart.js measurement histograms, statevector tables, and execution statistics.
    \item Achievement dashboard with XP bar, level display, and badge collection.
    \item Community library with recipe browsing, forking, and tag-based filtering.
    \item Full dark and light theme support.
\end{itemize}

\section{Testing and Validation}

\subsection{Testing Strategy}
The verification approach combines three layers: (1) backend regression tests covering the converter, simulator, database layer, and authentication flows; (2) targeted feature checks for Iterations 3 and 4; and (3) manual exploratory testing for visual behaviour (circuit rendering, gamification progression).

\subsection{Unit Testing}
Unit tests validate: conversion correctness for all supported gates across Cirq, Qiskit, and PennyLane; OpenQASM 3.0 output validity; QSim simulation correctness against known quantum state distributions; authentication and RBAC logic; and gamification achievement criteria.

\subsection{Test Coverage}

\begin{table}[H]
\centering
\caption{Test Coverage Summary}
\label{tab:test-coverage}
\begin{tabular}{|l|r|l|}
\hline
\textbf{Iteration / Suite} & \textbf{Tests} & \textbf{Scope} \\ \hline
Iteration 1 & 73 & Basic conversion, gates, type system \\ \hline
Iteration 2 & 151 & Advanced features, DB models, auth, RBAC \\ \hline
Integration Tests & 80 & End-to-end conversion and hybrid execution \\ \hline
\textbf{Total} & \textbf{304} & \textbf{100\% pass rate} \\ \hline
\end{tabular}
\end{table}

\textit{Note: Iterations 3 and 4 feature additions (gamification service, recipe sharing, SDK packaging) were validated via manual exploratory testing and integration test extensions included in the 80 integration tests above.}

\section{Evaluation and Benchmarking}

Iteration 4 delivered a reproducible benchmarking study titled \textit{``Cross-Framework Quantum Algorithm Benchmarking via Unified Compilation to OpenQASM 3.0''}, using QCanvas as the compilation platform and QSim as the execution engine.

\subsection{Benchmark Setup}
\begin{itemize}
    \item \textbf{Environment:} Windows 11, Intel Core i7, 16 GB RAM, Python 3.11.
    \item \textbf{Algorithm Suite:} 12 standard algorithms (Bell State, GHZ State, Quantum Teleportation, Deutsch--Jozsa, Bernstein--Vazirani, Grover's Algorithm, QRNG, VQE, QAOA, QML XOR Classifier, Bit-Flip Code, Phase-Flip Code).
    \item \textbf{Frameworks:} Qiskit, Cirq, PennyLane --- each with idiomatic native implementations.
    \item \textbf{Qubit Range:} 2--8 qubits per algorithm; 136 distinct configurations.
    \item \textbf{Shots:} 4,096 per simulation run; 10 runs per configuration for statistical robustness.
\end{itemize}

\subsection{Compilation Time Results}

Table~\ref{tab:compile-times} shows mean compilation latency (ms) for representative algorithms. All values are well within the NFR-P1 target of 1500ms.

\begin{table}[H]
\centering
\caption{Mean Compilation Latency (ms) -- Representative Algorithms at 8 Qubits}
\label{tab:compile-times}
\resizebox{\textwidth}{!}{%
\begin{tabular}{|l|r|r|r|}
\hline
\textbf{Algorithm} & \textbf{Qiskit (ms)} & \textbf{Cirq (ms)} & \textbf{PennyLane (ms)} \\ \hline
Bell State & 18.6 $\pm$ 0.4 & 8.3 $\pm$ 0.5 & 12.2 $\pm$ 0.3 \\ \hline
GHZ State & 21.9 $\pm$ 0.9 & 10.7 $\pm$ 0.4 & 38.0 $\pm$ 0.9 \\ \hline
Bernstein--Vazirani & 17.5 $\pm$ 0.6 & 13.0 $\pm$ 0.5 & 98.9 $\pm$ 2.5 \\ \hline
Deutsch--Jozsa & 7.0 $\pm$ 0.3 & 12.1 $\pm$ 1.3 & 106.3 $\pm$ 3.4 \\ \hline
Grover's Algorithm & 9.6 $\pm$ 0.5 & 32.5 $\pm$ 1.9 & 184.2 $\pm$ 6.5 \\ \hline
VQE & 32.9 $\pm$ 3.9 & 13.8 $\pm$ 1.0 & 55.2 $\pm$ 2.9 \\ \hline
QAOA & 55.0 $\pm$ 3.2 & 19.1 $\pm$ 2.6 & 73.7 $\pm$ 7.2 \\ \hline
\end{tabular}%
}
\end{table}

Cirq consistently achieves the lowest compilation latency (mean 8--35ms). PennyLane incurs the highest latency, exceeding 100ms for oracle-based algorithms due to its full gate decomposition strategy.

\subsection{Structural Benchmark Results}

\begin{table}[H]
\centering
\caption{Structural Benchmark Summary (Selected Algorithms)}
\label{tab:structural-summary}
\resizebox{\textwidth}{!}{%
\begin{tabular}{|l|l|r|r|r|r|}
\hline
\textbf{Algorithm} & \textbf{Framework} & \textbf{Qubits} & \textbf{Total Gates} & \textbf{Circuit Depth} & \textbf{Multi-Qubit Ratio} \\ \hline
Bell State & All Frameworks & 2 & 4 & 3 & 0.25 \\ \hline
\multirow{3}{*}{GHZ State} & Qiskit & 8 & 6 & 4 & 0.33 \\
 & Cirq & 8 & 6 & 4 & 0.33 \\
 & PennyLane & 8 & 16 & 9 & 0.44 \\ \hline
\multirow{3}{*}{Grover's Algorithm} & Qiskit & 6 & 18 & 12 & 0.11 \\
 & Cirq & 6 & 18 & 12 & 0.11 \\
 & PennyLane & 6 & 216 & 62 & 0.00 \\ \hline
\multirow{3}{*}{Bernstein--Vazirani} & Qiskit & 8 & 13 & 6 & 0.15 \\
 & Cirq & 8 & 13 & 6 & 0.15 \\
 & PennyLane & 8 & 31 & 4 & 0.16 \\ \hline
\multirow{3}{*}{VQE} & Qiskit & 6 & 5 & 2 & 0.20 \\
 & Cirq & 6 & 5 & 3 & 0.20 \\
 & PennyLane & 6 & 17 & 7 & 0.29 \\ \hline
\end{tabular}%
}
\end{table}

\begin{figure}[H]
    \centering
    \includegraphics[width=0.90\linewidth]{fig02_gate_counts.png}
    \caption{Gate Count Comparison Across Frameworks}
    \label{fig:gate-counts}
\end{figure}

\begin{figure}[H]
    \centering
    \includegraphics[width=0.90\linewidth]{fig03_circuit_depth.png}
    \caption{Circuit Depth Comparison Across Frameworks}
    \label{fig:circuit-depth}
\end{figure}

Qiskit and Cirq consistently produce structurally equivalent, compact circuits. PennyLane's default transpiler generates significantly more gates and deeper circuits for algorithms like Grover's (216 gates at 6 qubits vs.\ 18 for Qiskit/Cirq) due to its full-decomposition strategy for multi-qubit gates.

\subsection{Scaling Behaviour Analysis}

Power-law fits ($y = c \cdot n^a$) were computed for circuit depth as a function of qubit count. An exponent $a = 0$ indicates constant-depth scaling, optimal for NISQ execution.

\begin{table}[H]
\centering
\caption{Scaling Behaviour (Power-Law Fit: Circuit Depth vs.\ Qubits)}
\label{tab:scaling}
\begin{tabular}{|l|l|r|l|}
\hline
\textbf{Algorithm} & \textbf{Framework} & \textbf{Exponent ($a$)} & \textbf{Scaling Class} \\ \hline
GHZ State & Qiskit & 0.00 & Constant-depth \\
GHZ State & Cirq & 0.00 & Constant-depth \\
GHZ State & PennyLane & 1.00 & Linear \\ \hline
Deutsch--Jozsa & Qiskit & 0.00 & Constant-depth \\
Deutsch--Jozsa & Cirq & 0.00 & Constant-depth \\
Deutsch--Jozsa & PennyLane & 0.91 & Sub-linear \\ \hline
Grover's Algorithm & Qiskit & 0.00 & Constant-depth \\
Grover's Algorithm & Cirq & 0.00 & Constant-depth \\
Grover's Algorithm & PennyLane & 2.21 & Super-linear \\ \hline
Bernstein--Vazirani & Qiskit & 0.00 & Constant-depth \\
Bernstein--Vazirani & Cirq & 0.00 & Constant-depth \\
Bernstein--Vazirani & PennyLane & 0.87 & Sub-linear \\ \hline
\end{tabular}
\end{table}

\begin{figure}[H]
    \centering
    \includegraphics[width=0.80\linewidth]{fig05_ghz_scaling.png}
    \caption{GHZ State Depth Scaling Across Frameworks}
    \label{fig:ghz-scaling}
\end{figure}

\begin{figure}[H]
    \centering
    \includegraphics[width=0.80\linewidth]{fig14_grovers_scaling.png}
    \caption{Grover's Algorithm Depth Scaling --- PennyLane Super-Linear vs.\ Qiskit/Cirq Constant}
    \label{fig:grovers-scaling}
\end{figure}

\subsection{Statistical Equivalence Analysis}

Jensen--Shannon Divergence (JSD) was computed between pairwise output distributions to confirm that all three frameworks produce equivalent quantum computations despite structural differences.

\begin{table}[H]
\centering
\caption{Statistical Equivalence Summary (JSD, 4096 Shots)}
\label{tab:jsd}
\resizebox{\textwidth}{!}{%
\begin{tabular}{|l|r|r|r|r|r|r|}
\hline
\textbf{Algorithm} & \textbf{Qubits} & \textbf{JSD (Q--C)} & \textbf{JSD (Q--PL)} & \textbf{JSD (C--PL)} & \textbf{Q=C?} & \textbf{All Equiv.?} \\ \hline
Bell State & 2 & 0.0035 & 0.0002 & 0.0037 & \checkmark & \checkmark \\ \hline
GHZ State & 3 & 0.0005 & 0.0021 & 0.0017 & \checkmark & \checkmark \\ \hline
QML XOR Classifier & 2 & 0.0018 & 0.0049 & 0.0067 & \checkmark & \checkmark \\ \hline
QML XOR Classifier & 4 & 0.0010 & 0.0017 & 0.0007 & \checkmark & \checkmark \\ \hline
VQE & 2 & 0.0033 & 0.0029 & 0.0053 & \checkmark & \checkmark \\ \hline
QAOA & 2 & 0.0035 & 0.0035 & 0.0049 & \checkmark & \checkmark \\ \hline
Grover's Algorithm & 2 & 0.0000 & 0.0000 & 0.0000 & \checkmark & \checkmark \\ \hline
\end{tabular}%
}
\end{table}

\begin{figure}[H]
    \centering
    \includegraphics[width=0.85\linewidth]{fig09_jsd_heatmap.png}
    \caption{Jensen--Shannon Divergence Heatmap (Pairwise Framework Equivalence)}
    \label{fig:jsd-heatmap}
\end{figure}

Qiskit and Cirq are statistically equivalent (JSD $\approx 0$) across all tested algorithms, confirming that QCanvas's unified compilation pipeline preserves quantum semantics. PennyLane shows small divergence for multi-qubit oracle algorithms due to different transpilation strategies, but produces equivalent results for 2-qubit and variational circuits.

\subsection{Quantum Volume Analysis}

\begin{table}[H]
\centering
\caption{Effective Quantum Volume by Framework (Bell State Baseline)}
\label{tab:qv}
\begin{tabular}{|l|r|r|r|}
\hline
\textbf{Framework} & \textbf{$\log_2(\text{QV})$} & \textbf{QV} & \textbf{Fidelity (Bell)} \\ \hline
Qiskit   & 11 & 2048 & 0.9786 \\ \hline
Cirq     & 11 & 2048 & 0.9786 \\ \hline
PennyLane & 11 & 2048 & 0.9786 \\ \hline
\end{tabular}
\end{table}

All three frameworks achieve an effective QV of 2048 ($\log_2 = 11$) for ideal (noise-free) simulation, confirming that QSim provides a level simulation substrate. The divergence between frameworks in structural metrics (Section~\ref{tab:structural-summary}) is purely a consequence of differing compiler strategies, not of underlying quantum semantics.

\subsection{Holistic QPack Scores}

\begin{table}[H]
\centering
\caption{QPack Holistic Benchmark Scores}
\label{tab:qpack}
\begin{tabular}{|l|r|r|r|r|r|}
\hline
\textbf{Framework} & \textbf{$S_\text{compile}$} & \textbf{$S_\text{scale}$} & \textbf{$S_\text{accuracy}$} & \textbf{$S_\text{capacity}$} & \textbf{$S_\text{overall}$} \\ \hline
Qiskit   & 2.84 & 3.91 & 9.41 & 8 & \textbf{58.80} \\ \hline
Cirq     & 2.84 & 3.91 & 9.41 & 8 & \textbf{58.79} \\ \hline
PennyLane & 2.56 & $-$0.90 & 9.18 & 2 & 9.28 \\ \hline
\end{tabular}
\end{table}

\begin{figure}[H]
    \centering
    \includegraphics[width=0.80\linewidth]{fig15_qpack_radar.png}
    \caption{QPack Radar Chart --- Framework Holistic Score Comparison}
    \label{fig:qpack-radar}
\end{figure}

Qiskit and Cirq score near-identically (QPack $\approx 58.8$), substantially outperforming PennyLane (QPack $= 9.28$). PennyLane's negative scalability score ($S_\text{scale} = -0.90$) reflects its super-linear circuit growth for Grover's algorithm, which limits its applicability to larger NISQ problem sizes in default configuration.

\subsection{Performance Evaluation}

All measurements were taken on a Windows 11 host, Intel Core i7-12700H, 16 GB RAM, Python 3.11, with 10 repetitions per measurement and mean values reported.

\begin{itemize}
    \item \textbf{Conversion Latency:} Cirq circuits average 8--13ms; Qiskit 7--55ms; PennyLane 12--184ms depending on algorithm. All values satisfy NFR-P1 ($<$1500ms).
    \item \textbf{Simulation Throughput:} Grover's algorithm (6 qubits, 4096 shots) executes in under 800ms on the statevector backend (NFR-P2: $<$5 seconds).
    \item \textbf{Visualisation Render Time:} D3.js circuit diagram updates average sub-200ms for circuits up to 50 gates, satisfying NFR-P3.
    \item \textbf{Round-Trip Accuracy:} Converting standard gates to OpenQASM 3.0 and executing through QSim preserves measurement distributions with JSD $< 0.007$ against theoretical values.
\end{itemize}
